\Article{21}
Haben die Beobachtungen die (näherungsweise richtigen) Gleichungen \(v=0,v'=0,v''=0\) etc. geliefert, so ist zur Herleitung des Wertes der Unbekannten \(x\) aus ihnen eine Kombination dieser Gleichungen von der Form
\[
\kappa v+\kappa'v'+\kappa''v''+\text{etc.}=0
\]
zu verwenden, die \(x\) den Koeffizienten \(1\) gibt und die übrigen Unbekannten \(y,z\) etc. eliminiert; nach Art. 18 ist dieser Bestimmung das Gewicht
\[
=\frac1{\kappa\kappa+\kappa'\kappa'+\kappa''\kappa''+\text{etc.}}
\]
zuzuweisen. Aus dem vorigen Artikel folgt daher, dass die geeignetste Bestimmung diejenige sein wird, bei der \(\kappa=\alpha,\kappa'=\alpha',\kappa''=\alpha''\) etc. gesetzt wird. Auf diese Weise erhält \(x\) den Wert \(A\), und offensichtlich lässt sich derselbe Wert auch (ohne Kenntnis der Multiplikatoren \(\alpha,\alpha',\alpha''\) etc.) unmittelbar durch Elimination aus den Gleichungen \(\xi=0,\eta=0,\zeta=0\) etc. herleiten. Das dieser Bestimmung zuzuweisende Gewicht wird \(=\frac1{[\alpha\alpha]}\) sein oder der dabei zu befürchtende mittlere Fehler
\[
=m\BareRoot{p[\alpha\alpha]}=m'\BareRoot{p'[\alpha\alpha]}=m''\BareRoot{p''[\alpha\alpha]}\text{ etc.}
\]

Ganz ebenso wird die geeignetste Bestimmung der übrigen Unbekannten \(y,z\) etc. ihnen dieselben Werte geben, die durch Elimination aus denselben Gleichungen \(\xi=0,\eta=0,\zeta=0\) etc. hervorgehen.

Bezeichnen wir die unbestimmte Summe \(vv+v'v'+v''v''\text{ etc.}\) oder, was dasselbe ist, diese Summe
\[
p(V-L)^2+p'(V'-L')^2+p''(V''-L'')^2+\text{etc.}
\]
\SourceLeaf{030}mit \(\Omega\), so sind offenbar \(2\xi,2\eta,2\zeta\) etc. die partiellen Differentialquotienten der Funktion \(\Omega\), nämlich
\[
2\xi=\frac{\mathrm{d}\Omega}{\mathrm{d}x},\quad2\eta=\frac{\mathrm{d}\Omega}{\mathrm{d}y},\quad2\zeta=\frac{\mathrm{d}\Omega}{\mathrm{d}z}\text{ etc.}
\]
Die aus der geeignetsten Kombination der Beobachtungen hervorgehenden Werte der Unbekannten, die wir zweckmäßig die \textit{plausibelsten Werte} nennen können, werden daher mit denjenigen identisch sein, bei denen \(\Omega\) seinen kleinsten Wert erhält. Nun drückt \(V-L\) unbestimmt die Differenz zwischen dem berechneten und dem beobachteten Wert aus. Die plausibelsten Werte der Unbekannten werden also diejenigen sein, welche die Summe der mit den Beobachtungsgewichten multiplizierten Quadrate der Differenzen zwischen den beobachteten und berechneten Werten der Größen \(V,V',V''\) etc. zum Minimum machen. Diesen Grundsatz hatten wir in der \textit{Theoria Motus Corporum Coelestium} auf einem ganz anderen Weg begründet. Soll außerdem die relative Genauigkeit der einzelnen Bestimmungen angegeben werden, so müssen durch unbestimmte Elimination aus den Gleichungen (III) die Unbekannten \(x,y,z\) etc. selbst in folgender Form dargestellt werden:
\[
\left.\begin{aligned}
x&=A+[\alpha\alpha]\xi+[\alpha\beta]\eta+[\alpha\gamma]\zeta+\text{etc.}\\
y&=B+[\beta\alpha]\xi+[\beta\beta]\eta+[\beta\gamma]\zeta+\text{etc.}\\
z&=C+[\gamma\alpha]\xi+[\gamma\beta]\eta+[\gamma\gamma]\zeta+\text{etc.}
\end{aligned}\right\}\quad\text{(VII)}
\]
\begin{center}etc.\end{center}
Auf diese Weise werden die plausibelsten Werte der Unbekannten \(x,y,z\) etc. jeweils \(A,B,C\) etc. sein und die diesen Bestimmungen zuzuweisenden Gewichte \(\frac1{[\alpha\alpha]},\frac1{[\beta\beta]},\frac1{[\gamma\gamma]}\) etc. oder die dabei zu befürchtenden mittleren Fehler
\[
\begin{aligned}
\text{für }x\mathord{.}\,\mathord{.}\,\mathord{.}\,\mathord{.}\,\mathord{.}\,\mathord{.}\,&m\BareRoot{p[\alpha\alpha]}=m'\BareRoot{p'[\alpha\alpha]}=m''\BareRoot{p''[\alpha\alpha]}\text{ etc.}\\
\text{für }y\mathord{.}\,\mathord{.}\,\mathord{.}\,\mathord{.}\,\mathord{.}\,\mathord{.}\,&m\BareRoot{p[\beta\beta]}=m'\BareRoot{p'[\beta\beta]}=m''\BareRoot{p''[\beta\beta]}\text{ etc.}\\
\text{für }z\mathord{.}\,\mathord{.}\,\mathord{.}\,\mathord{.}\,\mathord{.}\,\mathord{.}\,&m\BareRoot{p[\gamma\gamma]}=m'\BareRoot{p'[\gamma\gamma]}=m''\BareRoot{p''[\gamma\gamma]}\text{ etc.}
\end{aligned}
\]
\begin{center}etc.\end{center}
in Übereinstimmung mit dem, was wir in der \textit{Theoria Motus Corporum Coelestium} gelehrt haben.

\Article{22}
Es empfiehlt sich, den einfachsten und zugleich häufigsten Fall, bei dem nur eine Unbekannte vorhanden ist und \(V=x,V'=x,V''=x\) etc. gilt, mit wenigen Worten gesondert zu behandeln. Es wird nämlich \(a=\BareRoot{p},a'=\BareRoot{p'},a''=\BareRoot{p''}\) etc., \(l=-L\BareRoot{p},l'=-L'\BareRoot{p'},l''=-L''\BareRoot{p''}\) etc. sein und folglich
\SourceLeaf{031}
\[
\xi=(p+p'+p''+\text{etc.})x-(pL+p'L'+p''L''+\text{etc.})
\]
Hieraus folgt weiter
\[
\begin{aligned}
\relax[\alpha\alpha]&=\frac1{p+p'+p''+\text{etc.}}\\[6pt]
A&=\frac{pL+p'L'+p''L''+\text{etc.}}{p+p'+p''+\text{etc.}}
\end{aligned}
\]
Ist also aus mehreren Beobachtungen von ungleicher Genauigkeit, deren Gewichte jeweils \(p,p',p''\) etc. sind, der Wert derselben Größe aus der ersten \(=L\), aus der zweiten \(=L'\), aus der dritten \(=L''\) etc. gefunden worden, so wird ihr plausibelster Wert
\[
=\frac{pL+p'L'+p''L''+\text{etc.}}{p+p'+p''+\text{etc.}}
\]
sein und das Gewicht dieser Bestimmung \(=p+p'+p''\text{ etc.}\) Haben alle Beobachtungen dieselbe Genauigkeit, so wird der plausibelste Wert
\[
=\frac{L+L'+L''+\text{etc.}}{\pi}
\]
sein, also gleich dem arithmetischen Mittel der beobachteten Werte, und das Gewicht dieser Bestimmung \(=\pi\), wenn das Gewicht der Beobachtungen als Einheit angenommen wird.
\par\nobreak\medskip\centerline{\rule{.25\linewidth}{.3pt}}
\clearpage
